% !TeX program = lualatex
% !TeX encoding = utf8
% !TeX spellcheck = uk_UA
% !TEX indexprogram = upmendex
% !TeX root =./LabWork_Part1.tex

%=========================================================
\chapter*{Оцінка похибок експериментальних результатів}
\label{ch:errors}
\addcontentsline{toc}{chapter}{Оцінка похибок експериментальних результатів}
%=========================================================

\section{Похибки вимірювань}
\label{sec:errors_intro}

Результат будь-якого вимірювання неідеальний: приладом можна отримати
або безпосередньо виміряну величину (\emph{пряме} вимірювання), або
величину, обчислену за формулою з кількох прямо виміряних величин
(\emph{непряме} вимірювання). У кожному із цих випадків отримане число
відрізняється від істинного значення $X$, тому результат вимірювання
без зазначення його точності не має фізичного сенсу.

\emph{Абсолютною похибкою} називають різницю
\begin{equation}\label{eq:delta}
	\delta X = X_{\text{вимір}} - X,
\end{equation}
а \emph{відносною похибкою} --- відношення
\begin{equation}\label{eq:reldelta}
	\varepsilon = \frac{\delta X}{X}.
\end{equation}
Оскільки істинне значення $X$ невідоме, ці похибки не можна обчислити
точно --- їх можна лише \emph{оцінити} за результатами вимірювань.

Похибки поділяють на два типи:
\begin{itemize}
	\item \textbf{статистичні (випадкові)} --- виникають через
	неконтрольовані впливи на прилад і на дослід, змінюються від
	вимірювання до вимірювання і можуть бути як додатними, так і
	від'ємними. Їх можна зменшити, збільшуючи кількість вимірювань, і
	оцінити методами математичної статистики;
	\item \textbf{систематичні} --- залишаються практично незмінними
	протягом усієї серії вимірювань (похибка приладу, неточність
	калібрування, недосконалість методики). Багаторазове повторення
	досліду їх не зменшує; єдиний спосіб боротьби з ними --- грамотне
	планування експерименту та поправки на відомі джерела зміщення.
\end{itemize}

Для більшості робіт цього практикуму (вимірювання кута закрутки
динамометра, показань електрометра, сили струму тощо) основний внесок
дає \emph{інструментальна} похибка приладу. Спосіб її оцінки залежить
від того, що вказано в паспорті приладу:
\begin{itemize}
	\item якщо прилад має лише шкалу без вказаного класу точності
	(наприклад, крутильний динамометр, лінійка, транспортир), похибку
	беруть за половину ціни поділки шкали;
	\item якщо в паспорті приладу вказано \emph{клас точності} $k$
	(наприклад, амперметр класу точності $0.5$ у Роботі~3), абсолютну
	похибку обчислюють за формулою
	\begin{equation}\label{eq:accuracy_class}
		\Delta X_\text{інст} = \frac{k}{100}\cdot X_\text{межа},
	\end{equation}
	де $X_\text{межа}$ --- границя вимірювання (кінцеве значення
	шкали, на якій встановлено прилад), а не поточне показання
	приладу. Через це відносна похибка вимірювання зростає, якщо
	стрілка приладу відхиляється лише на малу частку шкали, --- тому
	для точних вимірювань обирають границю, при якій показання
	приладу перевищує приблизно половину шкали.
\end{itemize}

\section{Похибки непрямих вимірювань}
\label{sec:propagation}

У жодній з робіт даного практикуму шукана величина ($q$, $F$, $M$, $k$,
$C$ тощо) не вимірюється безпосередньо --- вона обчислюється за
формулою через кілька прямо виміряних величин $x_1, x_2, \ldots, x_n$
з похибками $\Delta x_1, \Delta x_2, \ldots, \Delta x_n$. Якщо
\begin{equation}\label{eq:function}
	Y = f(x_1, x_2, \ldots, x_n),
\end{equation}
то абсолютну похибку $Y$ оцінюють за формулою поширення похибок:
\begin{equation}\label{eq:propagation}
	\Delta Y = \sqrt{\sum_{i=1}^{n} \left( \frac{\partial f}{\partial x_i} \right)^2 (\Delta x_i)^2}.
\end{equation}
Формула~\eqref{eq:propagation} справедлива у припущенні, що величини
$x_1, \ldots, x_n$ виміряні незалежно одна від одної (їхні похибки не
корельовані); для всіх робіт цього практикуму ця умова виконується.

На практиці зручніше користуватися логарифмічною похідною --- для
формул виду $Y = A \cdot x_1^{a_1} x_2^{a_2} \cdots x_n^{a_n}$
відносна похибка обчислюється простіше:
\begin{equation}\label{eq:rel_propagation}
	\frac{\Delta Y}{Y} = \sqrt{\sum_{i=1}^{n} a_i^2 \left( \frac{\Delta x_i}{x_i} \right)^2}.
\end{equation}

\textit{Приклад.} У Роботі~1 заряд кульки обчислюється за формулою
$q = 2.22\cdot 10^{-12}\cdot \varphi$, тому відносна похибка заряду
дорівнює відносній похибці встановленого потенціалу: $\Delta q / q =
\Delta \varphi / \varphi$. Натомість величина $A$ у степеневій
апроксимації $F = A\cdot (q^2)^{1+\sigma/2}$ містить уже кілька
прямо виміряних параметрів, і її похибку зручніше отримувати
безпосередньо з похибки коефіцієнтів апроксимації --- саме до цього
завдання ми переходимо в наступному розділі.

\section{Запис результату}
\label{sec:rounding}

Результат вимірювання величини $X$ записують у вигляді
\begin{equation}\label{eq:result}
	X = \left\langle X \right\rangle \pm \Delta X,
\end{equation}
при цьому необхідно правильно округлити обидва числа: зайві цифри
створюють оманливе враження точності, якої насправді немає.

Правила округлення:
\begin{enumerate}
	\item Округлення \textbf{починають з похибки}: якщо перша значуща
	цифра похибки --- $1$ або $2$, залишають \emph{дві} значущі цифри;
	якщо $3$ і більше --- залишають \emph{одну}.
	\begin{center}
		\begin{tabular}{cc}
			\toprule
			До округлення & Після округлення \\ \midrule
			$0.17295$     & $0.17$           \\
			$4.8329$      & $5$              \\
			$0.006298$    & $0.006$          \\ \bottomrule
		\end{tabular}
	\end{center}
	\item Далі округляють \textbf{саму величину} так, щоб її остання
	значуща цифра стояла на тій самій позиції, що й остання значуща
	цифра вже округленої похибки.
	\begin{center}
		\begin{tabular}{cc}
			\toprule
			До округлення      & Після округлення \\ \midrule
			$3.4874\pm0.17$    & $3.49\pm0.17$    \\
			$0.280184\pm0.006$ & $0.280\pm0.006$  \\ \bottomrule
		\end{tabular}
	\end{center}
	\item Якщо похибку записано з порядком $10^m$, той самий порядок
	вказують і для величини, взявши обидва числа в дужки:
	\begin{equation*}
		72903 \pm 400 \quad \longrightarrow \quad (72.9 \pm 0.4) \cdot 10^3.
	\end{equation*}
	Порядків $10^{-1}$ і $10^{1}$ бажано уникати, оскільки вони не дають
	«круглого» вигляду числа перед десятковою комою: для розмірних
	величин зручніше переходити до відповідної частинної одиниці (мілі-,
	кіло- тощо), а для безрозмірних --- записувати число без множника
	$10^m$ узагалі, якщо $0.1 \le X < 10$.
\end{enumerate}

\section{Графіки та метод найменших квадратів}
\label{sec:lsq}

У більшості робіт цього практикуму (Роботи~1, 3--5) результат
отримують саме з нахилу лінійної залежності виду
\begin{equation}\label{eq:linear}
	y = a x + b,
\end{equation}
побудованої за експериментальними точками. Правила побудови графіків:
\begin{itemize}
	\item графік супроводжують назвою, а осі --- позначенням величини
	та одиницею виміру;
	\item масштаб обирають так, щоб точки займали всю площу графіка;
	відлік не обов'язково починати з нуля, але це варто позначити
	розривом осі;
	\item кожну точку забезпечують \emph{хрестом похибок}, розміри
	якого відповідають $\Delta x$ і $\Delta y$ (не показують лише
	тоді, коли похибка настільки мала, що невидима на графіку);
	\item точки не з'єднують ламаною лінією; натомість наносять
	теоретичну (апроксимаційну) пряму;
	\item за наявності кількох серій даних для них використовують різні
	позначення й додають легенду.
\end{itemize}

Коефіцієнти $a$ (нахил) і $b$ (вільний член) прямої~\eqref{eq:linear}
знаходять методом найменших квадратів (МНК) за $N$ точками
$(x_i, y_i)$:
\begin{equation}\label{eq:lsq_ab}
	a = \frac{\overline{xy} - \bar{x}\,\bar{y}}{\overline{x^2} - \bar{x}^2}, \qquad
	b = \bar{y} - a\bar{x},
\end{equation}
де риска над величиною позначає середнє арифметичне за всіма точками
(наприклад, $\bar{x} = \frac{1}{N}\sum_i x_i$). Похибки коефіцієнтів
оцінюють за формулами
\begin{equation}\label{eq:lsq_errors}
	\Delta a = \sqrt{\frac{1}{N-2}\cdot\frac{\overline{y^2} - \bar{y}^2 - a^2(\overline{x^2}-\bar{x}^2)}{\overline{x^2}-\bar{x}^2}}, \qquad
	\Delta b = \Delta a \sqrt{\overline{x^2}}.
\end{equation}

Формули~\eqref{eq:lsq_ab}--\eqref{eq:lsq_errors} дають коректну оцінку
$a$, $b$ та їхніх похибок лише тоді, коли похибки всіх точок $y_i$
приблизно однакові, а похибками $x_i$ можна знехтувати: сама формула
для $\Delta a$ виводиться з розкиду точок відносно прямої (через
$\overline{y^2}$) і не використовує задані наперед $\Delta x_i$,
$\Delta y_i$ явно.

\section{Комп'ютерна обробка даних: \texttt{Python} та \texttt{Gnuplot}}
\label{sec:computer}

Навіть якщо умови застосовності формул~\eqref{eq:lsq_ab}--\eqref{eq:lsq_errors} виконуються, обчислювати
коефіцієнти $a$, $b$ та їхні похибки вручну для десятка
точок незручно. На практиці для
цього використовують спеціалізоване програмне забезпечення. Нижче
наведено мінімальні, але повністю робочі приклади лінійної
апроксимації в \texttt{Python} та \texttt{Gnuplot} для даних, поданих
у текстовому файлі \texttt{data.dat} у вигляді чотирьох стовпчиків
--- $x$, $y$ та їхні похибки $\Delta x$, $\Delta y$:

\bgroup
\begin{minted}[
	fontsize=\small,
	frame=lines,
	bgcolor=lightgray!20!white,
]{text}
x       y     dx      dy
0.022   0.4   0.005   0.1
0.038   0.8   0.005   0.1
0.058   1.2   0.005   0.1
0.090   1.6   0.005   0.1
0.101   2.0   0.005   0.1
0.123   2.4   0.005   0.1
0.130   2.8   0.005   0.1
\end{minted}
\egroup

\begin{Warning}
	\textbf{Типова помилка.} Не слід обчислювати коефіцієнт нахилу як
	середнє від частки $\left\langle y_i/x_i \right\rangle$ для кожної
	точки окремо. По-перше, усереднення допустиме лише для повторних
	вимірювань однієї й тієї самої величини, а не для точок різної
	прямої. По-друге, такий підхід не дозволяє перевірити саму лінійність
	залежності й ігнорує можливий вільний член $b$ (наприклад, через
	зсув нуля приладу). Правильний шлях --- побудувати графік $y(x)$,
	переконатися в його лінійності й провести \emph{єдину} найкращу
	пряму методом найменших квадратів~\eqref{eq:lsq_ab}.
\end{Warning}

\subsection{Обробка в \texttt{Python} (\texttt{scipy})}

Пакет \href{https://scipy.org}{\texttt{scipy}} містить модуль
\texttt{scipy.odr}, що реалізує метод ортогональної регресії
(\emph{Orthogonal Distance Regression}) --- на відміну від формул
\eqref{eq:lsq_ab}--\eqref{eq:lsq_errors}, які враховують похибку лише
по $y$, цей метод коректно враховує похибки \emph{обох} величин, $x$
і $y$ одночасно, що якраз і потрібно для більшості робіт цього
практикуму:

\bgroup
\begin{minted}[
	fontsize=\small,
	frame=lines,
	bgcolor=lightgray!20!white,
	breaklines,
]{python}
import numpy as np
from scipy.odr import ODR, Model, RealData
import matplotlib.pyplot as plt

# Завантаження даних: стовпчики x, y, dx, dy
x, y, dx, dy = np.loadtxt("data.dat", unpack=True)

# Модельна функція: y = a*x + b
def linear(beta, x):
    a, b = beta
    return a * x + b

model = Model(linear)
data = RealData(x, y, sx=dx, sy=dy)
odr = ODR(data, model, beta0=[1.0, 0.0])   # початкове наближення
output = odr.run()

a, b = output.beta
da, db = output.sd_beta
chi2_nu = output.res_var                   # приведений хі-квадрат

print(f"a = {a:.4g} +- {da:.4g}")
print(f"b = {b:.4g} +- {db:.4g}")
print(f"chi^2/(N-p) = {chi2_nu:.3g}")

# Графік
xx = np.linspace(x.min(), x.max(), 200)
plt.errorbar(x, y, xerr=dx, yerr=dy, fmt="o", label="Експеримент")
plt.plot(xx, linear([a, b], xx),
         label=f"$y = {a:.3g}\\,x + {b:.3g}$")
plt.xlabel("$x$")
plt.ylabel("$y$")
plt.legend()
plt.savefig("fit.png", dpi=150)
\end{minted}
\egroup

Величина $\chi^2_\nu$ (приведений хі-квадрат, поле \texttt{res\_var})
показує, наскільки добре модель узгоджується з даними з урахуванням
заявлених похибок: $\chi^2_\nu \approx 1$ означає добру відповідність;
$\chi^2_\nu \gg 1$ --- або модель неправильна, або похибки $\Delta
x_i$, $\Delta y_i$ занижені; $\chi^2_\nu \ll 1$ --- похибки, ймовірно,
завищені. Це значення варто наводити разом із результатом апроксимації.

Якщо похибками $x$ можна знехтувати (вони набагато менші за похибки
$y$), можна скористатися простішою функцією
\mintinline{python}{scipy.optimize.curve_fit} зі зваженням лише по
$y$ --- у цьому випадку результат практично збігається з
формулами~\eqref{eq:lsq_ab}--\eqref{eq:lsq_errors}:
\begin{minted}[fontsize=\small, frame=lines, bgcolor=lightgray!20!white, breaklines]{python}
from scipy.optimize import curve_fit

popt, pcov = curve_fit(lambda x, a, b: a*x + b,
                        x, y, sigma=dy, absolute_sigma=True)
a, b = popt
da, db = np.sqrt(np.diag(pcov))
\end{minted}

\subsection{Обробка в \texttt{Gnuplot}}

Той самий результат можна отримати значно коротшим скриптом у
\href{http://www.gnuplot.info}{\texttt{Gnuplot}}, вбудована команда
\texttt{fit} якого також мінімізує зважену суму квадратів відхилень:

\bgroup
\begin{minted}[
	fontsize=\small,
	frame=lines,
	bgcolor=lightgray!20!white,
	breaklines,
]{gnuplot}
# Модель
f(x) = a*x + b

# Апроксимація з вагами по y (за похибками у 4-му стовпчику)
set fit errorvariables
fit f(x) 'data.dat' using 1:2:4 yerror via a, b

print sprintf("a = %.4g +- %.4g", a, a_err)
print sprintf("b = %.4g +- %.4g", b, b_err)

# Графік з хрестами похибок і апроксимаційною прямою
set xlabel "x"
set ylabel "y"
plot 'data.dat' using 1:2:3:4 with xyerrorbars title 'Експеримент', \
     f(x) title sprintf('y = %.3g x + %.3g', a, b)
\end{minted}
\egroup

Скрипт запускають з командного рядка:
\begin{minted}[fontsize=\small, frame=lines, bgcolor=lightgray!20!white]{bash}
gnuplot -persist fit.gp
\end{minted}


\subsection{Результат апроксимації}

Застосувавши до наведеної вище таблиці даних метод ортогональної
регресії (\texttt{scipy.odr}, лістинг вище), отримуємо
\begin{equation}\label{eq:example_result}
	a = 20.7 \pm 1.3, \qquad b = -0.06 \pm 0.11,
\end{equation}
і --- після округлення за правилами розд.~\ref{sec:rounding} ---
апроксимаційну пряму
\begin{equation*}
    y = (20.7\pm1.3)\,x - (0.06\pm0.11).
\end{equation*}

Побудований графік цієї апроксимації
разом з експериментальними точками та хрестами похибок наведено на
рис.~\ref{fig:example_fit}.

\begin{figure}[h]
	\centering
	\begin{tikzpicture}
		\begin{axis}[
			width=0.75\textwidth, height=0.55\textwidth,
			xlabel={$x$},
			ylabel={$y$},
			ylabel style={xshift=-8pt},
			xmin=0, xmax=0.14,
			ymin=0, ymax=3.0,
			x tick label style={/pgf/number format/fixed},
			legend pos=north west,
			grid=major,
			grid style={gray!25},
			axis lines=left,
			tick align=outside,
		]
			% Апроксимаційна пряма y = a x + b (ODR: a = 20.661, b = -0.0588)
			\addplot[domain=0:0.14, samples=2, thick, red]
				{20.661*x - 0.0588};
			\addlegendentry{$y = 20.7\,x - 0.06$}

			% Експериментальні точки з хрестами похибок
			\addplot[
				only marks,
				mark=*,
				mark size=1.7pt,
				color=blue,
				error bars/.cd,
					x dir=both, x explicit,
					y dir=both, y explicit,
				error bar style={blue, thin},
			] table[x=x, y=y, x error=dx, y error=dy] {
				x      y     dx      dy
				0.022   0.4   0.005   0.1
				0.038   0.8   0.005   0.1
				0.058   1.2   0.005   0.1
				0.090   1.6   0.005   0.1
				0.101   2.0   0.005   0.1
				0.123   2.4   0.005   0.1
				0.130   2.8   0.005   0.1
			};
			\addlegendentry{Експеримент}
		\end{axis}
	\end{tikzpicture}
	\caption{Приклад лінійної апроксимації методом ортогональної
		регресії: експериментальні точки з хрестами похибок $\Delta x$,
		$\Delta y$ та апроксимаційна пряма~\eqref{eq:example_result}.}
	\label{fig:example_fit}
\end{figure}


\begin{Warning}
	\textbf{«Величина менша за свою похибку --- це помилка?»} Ні, і
	це одна з найчастіших причин розгубленості на захисті лабораторної
	роботи, тому розберемось детально.

	Похибка $\Delta b$ --- це не «діапазон, у якому обов'язково лежить
	істинне значення», а оцінка \emph{ширини інтервалу}, у межах якого
	з певною ймовірністю (для стандартної похибки --- приблизно
	$68\,\%$) могло б опинитися істинне значення параметра. Якщо
	записати результат як $b = -0.06 \pm 0.11$, це означає: істинне
	значення $b$ із помітною ймовірністю лежить десь у проміжку від
	$-0.17$ до $0.05$ --- і нуль \emph{потрапляє} в цей проміжок.

	Інакше кажучи: за наявною точністю вимірювань ми \textbf{не можемо
	статистично відрізнити} отримане $b$ від нуля. Це \emph{не} ознака
	зіпсованого експерименту чи помилки в обчисленнях --- навпаки, у
	даному прикладі це очікуваний і \emph{добрий} результат: теорія
	передбачає пряму пропорційності $y = ax$ (без вільного члена), а
	експеримент підтвердив, що додатковий зсув $b$, якщо він і є, надто
	малий, щоб його побачити при такій точності вимірювань.

	Як діяти в такому випадку:
	\begin{itemize}
		\item \textbf{Не} відкидайте $b$ і \textbf{не} «підганяйте»
		апроксимацію під $b = 0$ штучно (наприклад, вручну прибираючи
		вільний член із формули) --- наведіть результат апроксимації
		як є, з обома параметрами та їхніми похибками;
		\item у висновках прямо напишіть, що отримане значення $b$
		статистично узгоджується з нулем, і поясніть, чому це підтверджує
		(а не спростовує) теоретичну модель;
		\item якщо ж, навпаки, $|b|$ помітно \emph{більше} за $\Delta b$
		(наприклад, $b = 0.5 \pm 0.1$), це вже статистично значущий
		результат --- ненульовий вільний член, який варто пояснити
		фізично (наприклад, систематичним зсувом нуля приладу) чи
		вказати на межі застосовності моделі.
	\end{itemize}

	Той самий принцип стосується \emph{будь-якого} результату
	вимірювання, а не лише коефіцієнтів апроксимації: якщо
	$X = \left\langle X \right\rangle \pm \Delta X$ і
	$\left|\left\langle X
    \right\rangle\right|
    \le
    \Delta X$,
це
	означає лише, що точність вимірювання недостатня, аби впевнено
	стверджувати про ненульове значення $X$ --- і аж ніяк не те, що
	вимірювання «не вдалося».
\end{Warning}


Підсумовуючи: кінцевий результат щоразу оформлюють
за єдиним правилом --- у вигляді~\eqref{eq:result} з коректним
округленням (розд.~\ref{sec:rounding}), а похибку опосередкованих
величин, обчислених із коефіцієнтів апроксимації, знаходять за
формулою поширення похибок~\eqref{eq:rel_propagation} з
розд.~\ref{sec:propagation}.

\clearpage